\section{任务定义}

给定医学问题 $q$、生物医学语料库 $C$ 以及可选知识源 $K$，系统返回排序后的证据列表 $E$，并可选返回答案 $a$。主要任务是证据检索与重排序；答案生成是次要任务。

对于每个问题，第一阶段检索器从 $C$ 中返回候选池 $D_q = \{d_1,\ldots,d_M\}$。重排序器随后利用检索分数、生物医学语义分数、概念与实体特征，以及局部超图特征，对 $D_q$ 估计改进排序。评估重点包括 Recall@$k$、MRR@$k$、nDCG@$k$ 和 evidence coverage。

形式化输出为：
\[
  (a, E, P),
\]
其中 $a$ 在当前证据检索为主的设置中是可选项，$E$ 是排序后的支持证据列表，$P$ 是用于案例分析的可解释实体、概念或超图路径。

